\section{Directional Derivatives}

\begin{enumerate}
  \item What is the definition of the directional derivative of a function $f(x,y)$ at a point $(x_0, y_0)$ in the direction of a vector $\mathbf{v}$?
  \item Why must the direction vector be normalized when computing a directional derivative?
  \item Write the limit definition of the directional derivative.
  \item How is the directional derivative related to the gradient $\nabla f(x_0, y_0)$?
  \item State the formula for the directional derivative in terms of the gradient and a unit vector $\mathbf{u}$.
  \item What is the geometric interpretation of the directional derivative?
  \item Compute the directional derivative of $f(x,y) = x^2 + y^2$ at $(1,2)$ in the direction of $\mathbf{v} = (3,4)$.
  \item What is the direction of maximum increase of a function at a point?
  \item What is the maximum value of the directional derivative at a point, and how is it computed?
  \item In what direction is the directional derivative equal to zero?
  \item How are level curves of a function related to directional derivatives?
  \item Explain why the gradient is perpendicular to level curves.
  \item How can directional derivatives be used to test whether a function is constant along a curve?
  \item What is the relationship between the directional derivative and the total derivative?
  \item How do you compute the directional derivative for a function $f(x,y,z)$ in $\mathbb{R}^3$?
  \item Compute the directional derivative of $f(x,y,z) = x y z$ at $(1,1,1)$ in the direction of $\mathbf{v} = (1,1,1)$.
  \item What happens if the gradient at a point is the zero vector?
  \item Can a function have all directional derivatives at a point but still fail to be differentiable there?
  \item How is the directional derivative related to partial derivatives?
  \item Show how the directional derivative can be expressed as a dot product.
  \item What is the difference between directional derivatives and partial derivatives?
  \item How does the chain rule relate to directional derivatives?
  \item Given a parameterized curve $\mathbf{r}(t)$, how do you compute the rate of change of $f$ along the curve?
  \item Explain the concept of the derivative of $f$ along a path.
  \item How can directional derivatives be interpreted as slopes of tangent lines to curves on surfaces?
\end{enumerate}
